\documentclass[../../../include/open-logic-section]{subfiles}

\begin{document}

\olfileid{sth}{z}{union}
\olsection{في الاتحاد}

أثبتت~\olref[sth][z][sep]{prop:intersectionsexist} التقاطعات. وأما
الاتحادات الاعتباطية فتحتاج إلى مسلّمة أخرى:

\begin{axiom}[الاتحاد] لكل مجموعة~$\OLId{latin-looped}{A}$ توجد المجموعة
$\bigcup \OLId{latin-looped}{A} = \Setabs{\OLId{latin-ordinary}{x}}{(\exists \OLId{latin-ordinary}{b} \in \OLId{latin-looped}{A}) \OLId{latin-ordinary}{x} \in \OLId{latin-ordinary}{b}}$.
\[
	\forall \OLId{latin-looped}{A} \exists \OLId{latin-looped}{U} \forall \OLId{latin-ordinary}{x}(\OLId{latin-ordinary}{x} \in \OLId{latin-looped}{U} \liff (\exists \OLId{latin-ordinary}{b} \in \OLId{latin-looped}{A})\OLId{latin-ordinary}{x} \in \OLId{latin-ordinary}{b})
\]
\end{axiom}

ويسوغ التصور التراكمي-التكراري هذه المسلّمة. فلتكن~$\OLId{latin-looped}{A}$ مجموعة؛ فإن~$\OLId{latin-looped}{A}$
تكونت في مرحلة~$\OLId{latin-looped}{S}$ ما بموجب~\stageshier. وكل عنصر من~$\OLId{latin-looped}{A}$ قد تكوّن
\emph{قبل}~$\OLId{latin-looped}{S}$ بموجب~\stagesacc؛ وبالوجه نفسه يكون كل عنصر في عنصر من
$\OLId{latin-looped}{A}$ قد تكوّن قبل~$\OLId{latin-looped}{S}$. فتكون \emph{تلك} المجموعات كلها متاحة قبل~$\OLId{latin-looped}{S}$
ليُجمع شملها في مجموعة عند~$\OLId{latin-looped}{S}$، وهذه المجموعة هي~$\bigcup \OLId{latin-looped}{A}$.

\end{document}


